\documentclass[10pt,a4paper]{amsart}
\usepackage[margin=19mm]{geometry}
\usepackage{amsmath,amssymb,mathtools}
\usepackage[scr=rsfs,cal=euler]{mathalfa}
\usepackage{xfrac}
\usepackage[unicode,hidelinks,bookmarks=false]{hyperref}
\newcommand{\ie}{i.e.\ }
\newcommand{\cf}{cf.\ }
\newcommand{\resp}{respectively\ }
\newcommand{\Subsection}{\subsection}
\providecommand{\nobreakdash}{\nobreak\hbox{-}}
\setlength{\emergencystretch}{2em}
\multlinegap0pt
\usepackage{amssymb}
\usepackage{xcolor}
\newtheorem{lemma}{Lemma}
\newtheorem{theorem}{Theorem}
\title{Enumerating separable derangements}
\author{Robert Dougherty-Bliss, Alejandro B. Galv\'an, Michaela A. Polley, David Shuster}
\date{Source: arXiv:2608.27583v1 (CC BY 4.0)}
\begin{document}
\maketitle

\noindent\emph{Source and licence.} Enumerating separable derangements. Version arXiv:2608.27583v1 (CC BY 4.0), distributed by arXiv under the Creative Commons Attribution 4.0 International licence, http://creativecommons.org/licenses/by/4.0/, as recorded in the arXiv licence metadata for this exact version and shown on the abstract page https://arxiv.org/abs/2608.27583v1. Full licence notice: \texttt{licenses/arxiv-2608.27583-CC-BY-4.0.txt}.\par\smallskip

\noindent\textit{Editorial scope.} Counted unit: the article's theorem theorem (source0.tex lines 595--601). The article's complete original proof of it is delivered with it (source0.tex lines 603--626, one \texttt{proof} environment of 24 lines). No mathematics is written, repaired or re-derived: the statement and the proof are the article's own lines, in the article's own notation, and no retained line is substituted or re-flowed.  Retained uncounted as the source-local context the proof needs: lines 556--561.  Nothing was omitted from any retained span, so the package carries no omission record.  No retained line contains a citation, so the wrapper carries no bibliography and no bibliography entry.  No retained line contains a figure environment, an image-inclusion command or drawing code, so no image is delivered and the package declares no asset dependency.  Editorial formatting only, in addition: an AMS \texttt{amsart} preamble that restates the article's own declarations and macros used by the retained lines, namely the environments \texttt{lemma}, \texttt{theorem} on separate counters, together with the standard packages those lines need.  The retained environments print in this document as Lemma 1, Theorem 1 in order of appearance, whereas the article prints the same items with the numbering of its own full document.  The complete native source of the article as distributed in the arXiv e-print for this exact version is bundled unchanged as \texttt{sources/208725/source0.tex}.
\begin{lemma}
    \label{lemma:bound}
    For any integer $n \geq 1$ we have $b_{2n} \geq s_n^2$ and $b_{2n + 1} \geq
    s_n (s_{n + 1} - s_n)$. For $n \geq 3$ we have $b_n \leq s_n - 2 s_{n - 1} +
    s_{n - 2}$.
\end{lemma}

\begin{theorem}
    There exists a constant $\beta > 0$ such that
    \begin{equation*}
        \beta n^{-3/2} \leq \frac{b_n}{s_n} \leq 0.69
    \end{equation*}
    for sufficiently large $n$.
\end{theorem}

\begin{proof}
    By Lemma~\ref{lemma:bound}, we have
    \begin{equation*}
        \frac{b_{2n}}{s_{2n}} \geq \frac{s_n^2}{s_{2n}} \sim C \left(\frac{2}{n}\right)^{3/2}.
    \end{equation*}
    This implies $b_{2n} / s_{2n} = \Omega(n^{-3/2})$. Similarly,
    \begin{equation*}
        \frac{b_{2n + 1}}{s_{2n + 1}} \geq \frac{s_n(s_{n + 1} - s_n)}{s_{2n + 1}}
        \sim C \left( \frac{2}{n} \right)^{3/2} (1 - \lambda^{-1}).
    \end{equation*}
    This implies $b_{2n + 1} / s_{2n + 1} = \Omega(n^{-3/2})$. Taken together,
    these statements imply $b_n / s_n = \Omega(n^{-3/2})$.

    For the other direction, the upper bound part of Lemma~\ref{lemma:bound}
    implies
    \begin{equation*}
        \frac{b_n}{s_n} \leq 1 - 2\frac{s_{n - 1}}{s_n} + \frac{s_{n - 2}}{s_n}.
    \end{equation*}
    Letting $n \to \infty$ in the right-hand side produces
    \begin{equation*}
        1 - \frac{2}{\lambda} + \frac{1}{\lambda^2} = 0.68629\dots.
    \end{equation*}
    In particular, $b_n / s_n \leq 0.69$ for sufficiently large $n$.
\end{proof}

\end{document}
